% Paper 4, §4 — Probing: the supervised correlational readout
% Thesis: a probe is a (typically linear) classifier trained on frozen
% activations to predict an external property; high-accuracy probe is the
% (correlational) signal that the property is linearly represented. We
% trace four probe families — Alain & Bengio 2017 (linear logistic
% regression), Hewitt & Manning 2019 (structural probe), Burns 2022 (CCS,
% unsupervised), Marks & Tegmark 2023 (mass-mean) — and the layer-position
% story (surface early, syntax mid, semantic late-mid).

\section{Probing}
\label{sec:probing}

A \emph{probe} is the supervised correlational readout: a (typically
linear) classifier trained on \emph{frozen} activations of a pre-trained
transformer to predict an external property — truth, sentiment, syntactic
role, factuality, sycophancy. Where the lens of \S\ref{sec:logit-lens-family}
reuses the model's own unembedding $W_U$ to ask \emph{what would the
model predict at layer $\ell$}, a probe trains a fresh classifier on
$\mathbf{h}^{(\ell, t)}$ to ask a different question: \emph{is the
property $y$ linearly decodable from the activation at layer $\ell$,
position $t$?} The implicit claim that motivates the methodology is
sharper than the question: if a high-accuracy probe exists for property
$y$ in $\mathbf{h}^{(\ell, t)}$, then the model \emph{linearly
represents} $y$ at that layer–position. Whether the model also
\emph{uses} that representation is the correlation–causation gap that
\S\ref{sec:correlation-causation} takes up; this section sets up the
methodology, the canonical four probe families, and the layer-position
story that lets a probe-accuracy curve diagnose \emph{where} a property
is encoded.

\subsection{Setup: the linear probe}
\label{sec:probe-setup}

Let $\mathcal{M}$ be a frozen pre-trained transformer with $L$ layers and
hidden size $D$ (Paper~1 notation). For an input sequence $\mathbf{x}$,
write $\mathbf{h}^{(\ell, t)}(\mathbf{x}) \in \mathbb{R}^D$ for the
residual-stream activation at layer $\ell$, position $t$. A
\emph{linear probe} for a binary property $y$ is a classifier
\begin{equation}
  \hat y \;=\; g(\mathbf{h}^{(\ell, t)})
  \;=\; \sigma\!\left(\mathbf{w}^\top \mathbf{h}^{(\ell, t)} + b\right),
  \label{eq:linear-probe}
\end{equation}
with $\mathbf{w} \in \mathbb{R}^D$, $b \in \mathbb{R}$, and $\sigma$ the
sigmoid \citep{alain-bengio-2017}\footnote{KB excerpt:
\texttt{kb/excerpts/alain-bengio-2017.md\#abstract}.}. For a
multiclass property $y \in \{1, \dots, K\}$ the probe is a softmax
classifier with $\mathbf{W} \in \mathbb{R}^{K \times D}$. The probe is
fit on a labeled dataset
$\mathcal{D} = \{(\mathbf{x}_i, y_i)\}$ by minimizing a regularized
cross-entropy loss
\begin{equation}
  \mathcal{L}(\mathbf{w}, b)
  \;=\; \frac{1}{|\mathcal{D}|}\sum_{i}\mathcal{L}_{\mathrm{CE}}\!\left(g(\mathbf{h}^{(\ell, t)}(\mathbf{x}_i)),\, y_i\right)
  \;+\; \lambda \|\mathbf{w}\|_2^2,
  \label{eq:probe-loss}
\end{equation}
with $\mathcal{M}$'s parameters held fixed throughout. The layer $\ell$
and the token position $t$ are hyperparameters on equal footing with
$\lambda$. Probe accuracy on a held-out test split is the headline
diagnostic; per the control-task convention of
\citet{hewitt2019-control-tasks}, it must be reported alongside an
\emph{accuracy delta} against a randomly-initialized-model baseline so
that the probe's expressivity does not masquerade as the model's
representation. Tensor-shape note: a batched cache of activations at
layer $\ell$ is $\mathsf{B}\!\times\!\mathsf{S}\!\times\!D$; the probe
output is $\mathsf{B}\!\times\!\mathsf{S}\!\times\!1$ (binary) or
$\mathsf{B}\!\times\!\mathsf{S}\!\times\!K$ (multiclass).

\citet{alain-bengio-2017} introduced this framework on Inception~v3 and
ResNet-50 and reported that ``the linear separability of features
increase[s] monotonically along the depth of the
model''\footnote{KB excerpt: \texttt{kb/excerpts/alain-bengio-2017.md\#abstract}
(verbatim).}: a class-label probe at layer $\ell$ does at least as well
as the same probe at any earlier layer. The result anchored the
``hierarchical-abstraction'' picture that subsequent NLP probing work
inherits — though, as we will see in \S\ref{sec:probe-layer-position},
the LLM picture is non-monotonic and property-dependent in ways the
vision result does not anticipate.

\subsection{The structural probe: from binary to geometric}
\label{sec:structural-probe}

Linear logistic regression in Eq.~\eqref{eq:linear-probe} treats $y$ as a
single class label. Many properties of interest — syntactic dependency
trees, coreference chains, semantic role frames — are higher-arity:
they are \emph{relations between} tokens, not labels \emph{on} tokens.
\citet{hewitt2019-structural-probe} generalize the framework with a
\emph{structural probe}: a learned linear transform
$\mathbf{B} \in \mathbb{R}^{k \times D}$ such that, in $\mathbf{B}$-space,
the squared $L_2$ distance between two tokens' activations approximates
the parse-tree distance between those tokens
\citep{hewitt2019-structural-probe}\footnote{KB excerpt:
\texttt{kb/excerpts/hewitt2019-structural-probe.md\#abstract}.}:
\begin{equation}
  \boxed{\;
  d_{\mathbf{B}}(\mathbf{h}_i, \mathbf{h}_j)^2
  \;:=\; \|\mathbf{B}(\mathbf{h}_i - \mathbf{h}_j)\|_2^2
  \;\approx\; d_T(w_i, w_j)
  \;}
  \label{eq:structural-probe}
\end{equation}
where $d_T(w_i, w_j)$ is the number of edges on the parse-tree path
between tokens $w_i$ and $w_j$. A companion variant probes scalar
\emph{tree depth} via $\|\mathbf{B}\mathbf{h}_i\|_2^2 \approx
\mathrm{depth}_T(w_i)$. The rank $k$ is a hyperparameter; the entire
probe is a single $k \times D$ matrix per (model, layer). Hewitt and
Manning report that such $\mathbf{B}$ exists for ELMo and BERT but
\emph{not} for randomly-initialized baselines, providing evidence that
``entire syntax trees are embedded implicitly in deep models' vector
geometry''~\citep{hewitt2019-structural-probe}. The methodological
point: a structural probe shifts the question from ``does the model
represent $y$?'' to ``does the model represent the \emph{structure} of
$y$?'', and the answer is geometric rather than classificatory.

\subsection{Mass-mean probing: closed-form, parameter-free}
\label{sec:mass-mean}

The 2023+ truth-direction literature found that a probe with no learned
parameters at all often suffices, and is in fact \emph{more} causally
implicated than its trained counterparts. Given a balanced labeled set
$\{(\mathbf{h}_i, y_i \in \{+, -\})\}$ at a chosen $(\ell, t)$, the
\emph{mass-mean} (or \emph{difference-in-means}) probe direction is
\citep{marks-tegmark-2023-truth}\footnote{KB excerpt:
\texttt{kb/excerpts/marks-tegmark-2023-truth.md\#abstract},
\texttt{\#sec-1-contributions}.}:
\begin{equation}
  \boxed{\;
  \mathbf{d}_{\mathrm{MM}}
  \;:=\; \boldsymbol{\mu}_+ - \boldsymbol{\mu}_-
  \;=\; \mathbb{E}_{i:\,y_i = +}[\mathbf{h}_i] \;-\; \mathbb{E}_{i:\,y_i = -}[\mathbf{h}_i]
  \;}
  \label{eq:mass-mean}
\end{equation}
The classifier is the projection $\mathbf{h}^\top \mathbf{d}_{\mathrm{MM}}$
thresholded at the midpoint
$\tfrac{1}{2}(\boldsymbol{\mu}_+ + \boldsymbol{\mu}_-)^\top
\mathbf{d}_{\mathrm{MM}}$. There is no optimization, no regularization,
no hyperparameter tuning beyond $(\ell, t)$, and no stochasticity beyond
the sample split. \citet{marks-tegmark-2023-truth} argue that despite
this simplicity, mass-mean ``performs as well as other probing techniques
while identifying directions which are more causally implicated in model
outputs''~\citep[abstract]{marks-tegmark-2023-truth}: the absence of
moving parts is the \emph{point}, since the discovered direction is then
a property of the data, not a property of an optimizer. The strong
form of the methodological claim — that the mass-mean direction is the
direction the model actually \emph{uses} — is the load-bearing
contribution of \S\ref{sec:correlation-causation}, and the geometry-of-truth
case study of \S\ref{sec:truth-directions} unpacks the evidence and the
caveats.

\subsection{Burns CCS: probing without labels}
\label{sec:ccs}

A fourth probe family drops the label requirement entirely. Contrast-Consistent
Search (CCS) \deepencite{burns2022-ccs §3, citation-pending} fits a
linear probe by minimizing a self-consistency loss across paired
contrast prompts (e.g., ``X is true.'' / ``X is false.''), enforcing
that the probe's outputs on the two completions of a contrast pair
should sum to one and be confident. The headline claim is that CCS
recovers a \emph{truth-like} direction in LLM activations without
ever observing a truth label, and that the recovered direction's
classification accuracy is competitive with supervised probes on
in-distribution examples. \citet{marks-tegmark-2023-truth} situate
CCS alongside logistic-regression and contrastive probes
\citep[\S 1.1]{marks-tegmark-2023-truth}\footnote{KB excerpt:
\texttt{kb/excerpts/marks-tegmark-2023-truth.md\#sec-1-related}.} and
report that subsequent work has documented generalization failures
under negation — CCS-discovered directions on `cities` do not transfer
cleanly to `neg\_cities`. The unsupervised commitment is methodologically
appealing (it removes the curated-label confound) but inherits a strong
form of the spurious-correlate failure: any consistent latent feature
that distinguishes the two halves of the contrast set, not just truth,
is a candidate. The deeper treatment of CCS's geometry — why it works,
where it fails, what it shares with mass-mean — is in
\S\ref{sec:truth-directions}.

\begin{intuition}
A probe is a \emph{point-spread function} on the model's representation
space. Like the optical PSF of a microscope, it is the linear filter
$\mathbf{w}^\top \mathbf{h}$ in Eq.~\eqref{eq:linear-probe} through which
we observe the model's activation at $(\ell, t)$ — and the property $y$
must be \emph{in focus along that filter axis} for the probe to detect
it. The mass-mean direction
$\mathbf{d}_{\mathrm{MM}} = \boldsymbol{\mu}_+ - \boldsymbol{\mu}_-$ in
Eq.~\eqref{eq:mass-mean} is a closed-form choice of PSF: the axis along
which the two classes are maximally separated in expectation. The
structural probe in Eq.~\eqref{eq:structural-probe} replaces the scalar
projection with a $k$-dimensional one, so the PSF becomes a
$k$-dimensional subspace rather than a single axis. Returning to
canonical form: in every case the probe is a linear map
$\mathbb{R}^D \to \mathbb{R}^k$ (with $k = 1$ for binary, $k > 1$ for
structural) followed by a fixed decision rule, and probe accuracy
quantifies how well the chosen $k$-dimensional subspace separates the
labeled data.
\end{intuition}

\subsection{Where to probe: layer and token position}
\label{sec:probe-layer-position}

Probe accuracy is a function of $(\ell, t)$, and the shape of that
function tells a structural story about \emph{where} a property is
represented. Three regimes recur across the LLM probing literature.

\paragraph{Surface features early.} Token identity, basic part-of-speech
role, and other surface-syntactic properties probe highest in
\emph{early} layers; the probe accuracy curve rises steeply through the
embedding and first few transformer blocks and then plateaus or decays
\citep{hewitt2019-structural-probe}. The structural probe of
\S\ref{sec:structural-probe} is the canonical demonstration: parse-tree
distances are recovered with $\mathbf{B}$ trained on early-to-mid BERT
layers.

\paragraph{Syntax mid.} Tree-structure properties — dependency-parse
distances, constituent boundaries, agreement features — peak in
\emph{middle} layers. \citet{alain-bengio-2017} report monotonic
linear separability with depth on vision class labels, but the LLM
generalization is non-monotonic: information about syntactic structure
is most accessible mid-stack and degrades thereafter as later layers
recombine it for downstream prediction.

\paragraph{Semantic and pragmatic features late-mid.} Truth, factuality,
sentiment, speaker stance — properties that require the model to
integrate information across a clause — peak in \emph{late-mid} layers
and stay accurate through the end of the stack. The
\citet{marks-tegmark-2023-truth} finding on LLaMA-2-13B is the canonical
worked example: the truth direction localizes around layers 12--15 and
remains salient through layer 35, with PCA visualizations showing the
linear true/false separation emerging in early-middle layers and
``emerg[ing] later for datasets of more structurally complex statements
(e.g.\ conjunctive statements)''~\citep[\S 4]{marks-tegmark-2023-truth}\footnote{KB
excerpt: \texttt{kb/excerpts/marks-tegmark-2023-truth.md\#sec-3-layers},
\texttt{\#sec-4-misalignment}.}.

\paragraph{Position selection: the period-token trick.} For free-form
text, the most-informative probe position is often the
\emph{end-of-statement punctuation token} (typically the period). The
intuition is summarization: the period sits one position past the last
content token, and the model's attention pattern at that position has had
the opportunity to summarize the preceding clause into a single
hidden-state vector. \citet[\S 3]{marks-tegmark-2023-truth} document this
on LLaMA-2 family: causal-tracing localizes truth representations to a
group of hidden states ``over the final token of the statement and
end-of-sentence punctuation,'' which they hypothesize ``store a
representation of the statement's truth''\footnote{KB excerpt:
\texttt{kb/excerpts/marks-tegmark-2023-truth.md\#sec-3-layers}.}. The
period-token convention is now standard in the truth-direction
literature; \citet{marks-tegmark-2023-truth} cite Tigges et al.\ 2023 as
the prior empirical observation of the same summarization behavior on
sentiment.

\subsection{Probe lineage: eight families on one timeline}
\label{sec:probe-lineage}

\begin{table}[t]
\centering
\caption{The eight probe families that anchor the LLM probing literature.
Compiled from \texttt{kb/notes/interpretability/probing.md} \S 4.1.}
\label{tab:probe-families}
\small
\begin{tabularx}{\linewidth}{l l X X}
\toprule
\textbf{Family} & \textbf{Year} & \textbf{Probe form} & \textbf{Key limitation} \\
\midrule
Linear logistic regression \citep{alain-bengio-2017}
  & 2017
  & $\sigma(\mathbf{w}^\top \mathbf{h} + b)$ on frozen $\mathbf{h}$
  & correlational; needs $\ell_2$ regularization to avoid overfit \\
MLP probe
  & 2010s
  & small feedforward net on $\mathbf{h}$
  & probe leakage: a powerful probe learns $y$ from low-information features \citep{hewitt2019-control-tasks} \\
Structural probe \citep{hewitt2019-structural-probe}
  & 2019
  & $\|\mathbf{B}(\mathbf{h}_i - \mathbf{h}_j)\|^2 \approx d_T(w_i, w_j)$
  & confounded with task-difficulty when comparing layers \\
Control-task baseline \citep{hewitt2019-control-tasks}
  & 2019
  & same probe on randomly-permuted labels
  & required ablation; inflates accuracy delta on poorly-initialized models \\
Logistic regression for truth \citep{marks-tegmark-2023-truth}
  & 2023
  & linear, $\ell_2$-regularized, on truth labels
  & correlational unless paired with ablation \\
Mass-mean / difference-in-means \citep{marks-tegmark-2023-truth}
  & 2023
  & $\boldsymbol{\mu}_+ - \boldsymbol{\mu}_-$, closed-form
  & requires balanced classes; binary only \\
Difference-in-means $+$ steering
  & 2025
  & mass-mean direction, applied as activation-steering
  & within-model only; cross-family transfer untested \\
Contrastive / CCS unsupervised
  & 2022
  & linear probe, self-consistency loss, no labels
  & generalization failures under negation \citep[\S 1.1]{marks-tegmark-2023-truth} \\
\bottomrule
\end{tabularx}
\end{table}

Table~\ref{tab:probe-families} traces the eight families that anchor the
LLM probing literature, in the form lifted from
\texttt{kb/notes/interpretability/probing.md} \S 4.1. The lineage tells
a coherent story. The 2017--2019 wave (rows 1--4) establishes the
methodology and the first methodology critique: the control-task
baseline of \citet{hewitt2019-control-tasks} is the explicit recognition
that probe accuracy alone is insufficient, since a sufficiently
expressive probe can learn $y$ from features the model itself does not
use. The 2022 review of \citet{belinkov2022-probing-survey}\footnote{KB
excerpt: \texttt{kb/excerpts/belinkov2022-probing-survey.md\#abstract}.}
consolidates the critique into three confounders — probe expressivity,
spurious correlation, vestigial information — and surveys the
``advances'' that pair probes with causal interventions
\citep[\S abstract; advances]{belinkov2022-probing-survey}. Rows 5--8 are
the 2022+ post-critique wave: closed-form probes that minimize moving
parts (mass-mean), unsupervised probes that remove the curated-label
confound (CCS), and the activation-steering family that turns a probe
direction into a behavioral intervention. \S\ref{sec:correlation-causation}
takes up the methodological turn that connects rows 1--4 and rows 5--8;
\S\ref{sec:truth-directions} unpacks the row-5--6 case study at the level
of cross-dataset alignment.

\subsection{The probing protocol in six steps}
\label{sec:probe-protocol}

Synthesizing across the four probe families, the practical recipe for
probing a property $y$ at layer $\ell$, position $t$ is the following
six-step pipeline:
\begin{enumerate}
  \item \textbf{Curate} a labeled dataset
    $\mathcal{D} = \{(\mathbf{x}_i, y_i)\}$ in \emph{minimal-pair} form:
    inputs that vary only on $y$ and on as few confounders as possible.
    \citet[\S 2]{marks-tegmark-2023-truth} construct
    \texttt{cities}, \texttt{neg\_cities}, \texttt{larger\_than},
    \texttt{cities\_cities\_conj}, and the \texttt{likely} control
    dataset (10\,000 nonfactual completions) precisely to control for
    surface confounders such as text probability\footnote{KB excerpt:
    \texttt{kb/excerpts/marks-tegmark-2023-truth.md\#sec-2-datasets}.}.
  \item \textbf{Cache} activations $\{\mathbf{h}^{(\ell, t)}_i\}_i$ via
    one forward pass per example. Cost is $|\mathcal{D}|$ forward
    passes plus storage for $|\mathcal{D}| \cdot D$ floats per
    $(\ell, t)$ slice.
  \item \textbf{Fit} a probe — Eq.~\eqref{eq:linear-probe} via logistic
    regression, Eq.~\eqref{eq:mass-mean} closed-form, or
    Eq.~\eqref{eq:structural-probe} via least-squares regression — on a
    train split. Multi-$(\ell, t)$ runs are common.
  \item \textbf{Evaluate} held-out accuracy on a same-distribution test
    split, paired with a control-task accuracy from
    \citet{hewitt2019-control-tasks} on randomly-permuted labels.
  \item \textbf{Generalization tests} (optional but increasingly
    standard). Train on $\mathcal{D}_1$, evaluate on $\mathcal{D}_2$ at
    a structurally distinct distribution; failures here reveal that
    the probe was reading a surface confounder rather than the abstract
    property. \citet[\S 1, \S 4]{marks-tegmark-2023-truth} make
    cross-dataset transfer the cornerstone of their analysis.
  \item \textbf{Causal validation} (optional, but the \emph{closing}
    step for any claim that the model \emph{uses} the probe direction).
    Ablate $\mathbf{d}_y$ from the residual stream during a forward pass
    and measure behavioral change; or add $\alpha \hat{\mathbf{d}}_y$
    to steer the model. The ablation operator and the steering rule
    are spelled out in \S\ref{sec:correlation-causation}.
\end{enumerate}

Steps 1--4 produce a correlational claim of the form ``$y$ is linearly
decodable from $\mathbf{h}^{(\ell, t)}$ with accuracy $a$ above a
control-task baseline''. Step~5 sharpens the correlational claim by
testing whether the decoded direction generalizes across surface
distributions. Step~6 converts a correlational claim into a causal one,
and is the methodological centerpiece of the next section.

\subsection{Cross-paper thread and forward link}
\label{sec:probe-forward}

Probes are read from \emph{training-time activations}: the property
$y$ is whatever the model has acquired by the end of pre-training (or
post-training, depending on which checkpoint the activations are cached
from), and the probe-accuracy curve is a snapshot of the model's
representational competence at that checkpoint. The training pipeline
of Paper~2 is therefore the upstream substrate of every claim made in
this section: probe-accuracy results read off a checkpoint at
post-pre-training time inherit whatever pre-training data, mixture, and
optimizer choices produced that checkpoint, and a probe-accuracy
\emph{shift} between two checkpoints (e.g., pre-SFT vs.\ post-SFT) is
itself a probing measurement of how a training stage moves the model's
internal representation of $y$.

The methodological gap that probing alone leaves open — that a
high-accuracy probe is correlational, not causal — is the
load-bearing critique that
\S\ref{sec:correlation-causation} closes by pairing probes with
activation interventions: direction ablation
$\mathbf{h} \mapsto \mathbf{h} - (\mathbf{h}^\top \hat{\mathbf{d}}_y)\hat{\mathbf{d}}_y$
tests whether the direction is causally used, and direction addition
$\mathbf{h} \mapsto \mathbf{h} + \alpha \hat{\mathbf{d}}_y$ is the
activation-steering family that turns a probe direction into an
inference-time control surface. \S\ref{sec:truth-directions} then
returns to the geometry-of-truth case study at full depth: when probes
across datasets agree at late layers but diverge at early layers, what
does the divergence tell us about whether ``truth'' is a single linear
axis or a multi-axis subspace?
